% META: {"id": "TSPE-LIMFCT-EX-021", "chapitre": "TSPE-LIMITES-FONCTIONS", "type_objet": "exercice", "capacites": ["TSPE-LIMFCT-C2"], "capacites_codes": ["C2"], "methodes": ["M2"], "parcours": 1, "competences": ["calculer"], "duree_min": 8, "mode_creation": "ex_nihilo", "sources_inspiration": [], "fichier_tex": "chapitres/TSPE-LIMITES-FONCTIONS/exercices/TSPE-LIMFCT-EX-021.tex", "corrige_tex": "chapitres/TSPE-LIMITES-FONCTIONS/corriges/TSPE-LIMFCT-CO-021.tex", "status": "generated"}
% BEGIN-VERIFY
% from sympy import *
% x = symbols('x')
% f = (3*x + 1)/(x - 2)
% assert limit(f, x, oo) == 3
% assert limit(f, x, 2, '+') == oo
% assert limit(f, x, 2, '-') == -oo
% END-VERIFY
\begin{exercice}{TSPE-LIMFCT-EX-021}{1}{8}
Soit $f(x) = \dfrac{3x+1}{x-2}$ definie sur $\mathbb{R} \setminus \{2\}$.
\begin{enumerate}
  \item Determiner les limites de $f$ en $+\infty$ et en $-\infty$.
  \item Determiner les limites de $f$ en $2^+$ et $2^-$.
  \item En deduire les asymptotes de $\mathcal{C}_f$.
\end{enumerate}
\end{exercice}
